
\documentclass[11pt,a4paper]{article}
\usepackage[utf8]{inputenc}
\usepackage[T1]{fontenc}
\usepackage[spanish,es-tabla]{babel}
\usepackage{lmodern}
\usepackage[a4paper,left=1.8cm,right=1.8cm,top=2.15cm,bottom=2.0cm]{geometry}
\usepackage{microtype}
\usepackage[table]{xcolor}
\usepackage{array,longtable,booktabs,tabularx,multirow}
\usepackage{enumitem}
\usepackage{amsmath}
\usepackage{fancyhdr}
\usepackage{lastpage}
\usepackage{hyperref}
\usepackage{url}
\usepackage{ragged2e}
\usepackage{needspace}
\hypersetup{colorlinks=true,linkcolor=black,urlcolor=blue!50!black,citecolor=blue!45!black,pdftitle={FORM-TL-01 - Taller de tipos de requisitos},pdfauthor={Estudiante de Ingenieria de Software}}

\definecolor{verdeUTEQ}{HTML}{087647}
\definecolor{verdeTenue}{HTML}{EDF6F0}
\definecolor{grisTenue}{HTML}{F3F5F6}
\definecolor{grisTexto}{HTML}{404850}
\newcommand{\NombreEstudiante}{Chavarria Cuenca Tahiny Mel}
\newcommand{\RepoTexto}{https://github.com/TahinyMel/ISR401-ChavarriaTahiny/tree/main/FORM-TL-01}
\newcommand{\FechaRevision}{5 de noviembre de 2026}
\newcommand{\Docente}{Gleiston Cicerón Guerrero Ulloa, PhD}
\newcommand{\SiNoAplica}{\textit{No aplica (proceso)}}
\newcommand{\Tipo}[1]{\textbf{#1}}
\setlength{\parskip}{5pt}
\setlength{\parindent}{0pt}
\setlist[itemize]{leftmargin=*,nosep,topsep=2pt}
\setlist[enumerate]{leftmargin=*,itemsep=3pt}
\renewcommand{\arraystretch}{1.2}
\pagestyle{fancy}
\fancyhf{}
\fancyhead[L]{\footnotesize\color{grisTexto}UTEQ · Ingeniería de Requerimientos · ISR-401}
\fancyhead[R]{\footnotesize\color{grisTexto}FORM-TL-01}
\fancyfoot[L]{\footnotesize\color{grisTexto}Taller de tipos de requisitos}
\fancyfoot[R]{\footnotesize\color{grisTexto}\thepage\ / \pageref{LastPage}}
\renewcommand{\headrulewidth}{0.35pt}
\renewcommand{\footrulewidth}{0pt}
\setlength{\headheight}{14pt}
\newcommand{\encabezadoTabla}{\rowcolor{verdeUTEQ}\color{white}}
\newcolumntype{P}[1]{>{\RaggedRight\arraybackslash}p{#1}}

\begin{document}
\begin{center}
{\Large\bfseries\color{verdeUTEQ} UNIVERSIDAD TÉCNICA ESTATAL DE QUEVEDO}\\[4pt]
{\small Facultad de Ciencias de la Computación · Carrera de Software}\\[13pt]
{\large\bfseries Ingeniería de Requerimientos (ISR-401)}\\[5pt]
{\LARGE\bfseries Taller de tipos de requisitos}\\[4pt]
{\large FORM-TL-01}\\[9pt]
{\small\itshape Estudio de caso: Plataforma de gestión de citas médicas\\
Policlínico Ribera del Quevedo (organización ficticia)}
\end{center}
\vspace{0.3cm}
\begin{tabularx}{\textwidth}{>{\bfseries}p{3.4cm}X}
Estudiante: & \NombreEstudiante\\
Docente: & \Docente\\
Nivel y paralelo: & Cuarto nivel, paralelo B\\
Modalidad: & Individual\\
Fecha de revisión: & \FechaRevision\\
Fecha límite: & 7 de noviembre de 2026, 23:59\\
Repositorio: & \RepoTexto\\
\end{tabularx}
\vspace{0.25cm}


\section{Objetivo y delimitación del caso}
Clasificar los 20 enunciados del caso en tres dimensiones: (i) funcional, calidad o restricción; (ii) producto o proceso; y (iii) sistema o software. Para los enunciados de calidad se identifica la característica y subcaracterística pertinente del modelo ISO/IEC 25010:2023. Se transforman cinco expresiones vagas en requisitos verificables con métrica, unidad y umbral, y se deja preparada la comparación con un compañero, como solicita la guía \cite{guia2026,pohl2025,iso29148,iso25010}.

El caso describe un policlínico ficticio con 24 médicos, 9 especialidades, 18 consultorios, aproximadamente 12\,000 pacientes y 450 citas diarias. Cerca del 35\,\% de sus pacientes supera los 60 años. El servicio opera de lunes a sábado, de 07:00 a 19:00, y el pico de usuarios simultáneos esperado es de unos 300, especialmente los lunes a las 07:00. El sistema de facturación ya existe y proporciona servicios web; no se reemplazará. Tampoco forman parte del alcance la historia clínica digital, el laboratorio, la farmacia ni la telemedicina \cite{guia2026}.

\section{Criterios aplicados en la clasificación}
\begin{itemize}
  \item \textbf{Funcional:} prescribe una capacidad o comportamiento observable del sistema ante un evento o usuario. Una regla de negocio automatizada también se considera funcional.
  \item \textbf{Calidad:} especifica una propiedad medible del producto, como tiempo de respuesta, confidencialidad, inclusividad, disponibilidad, interoperabilidad o recuperación. Se usa el modelo de nueve características de ISO/IEC 25010:2023 \cite{iso25010}.
  \item \textbf{Restricción:} delimita alternativas de diseño, plataformas, normativa o condiciones de ejecución del proyecto.
  \item \textbf{Producto / proceso:} ``producto'' describe lo que ha de presentar o cumplir la solución entregada; ``proceso'' regula cómo se trabaja, valida o entrega el proyecto.
  \item \textbf{Sistema / software:} ``sistema'' incluye la solución sociotécnica completa, equipamiento, infraestructura e interfaces externas; ``software'' se refiere principalmente a una aplicación o comportamiento programado. En requisitos \emph{puramente de proceso} se marca \textit{no aplica}, ya que no definen el nivel de la solución; se indica el ámbito del proceso en la justificación. Esta es una decisión de clasificación explícita, susceptible de discusión en clase.
\end{itemize}

\section{Los 20 enunciados originales}
Los siguientes enunciados proceden del apartado 1.7 del caso (anexo de la guía); se preserva su sentido literal y su código de trazabilidad \cite{guia2026}.
{\small
\begin{longtable}{@{}P{1.4cm}P{15.55cm}@{}}
\caption{Enunciados E01--E20 del caso}\label{tab:enunciados}\\
\toprule \rowcolor{verdeUTEQ}\color{white}\textbf{Código} & \color{white}\textbf{Enunciado}\\\midrule\endfirsthead
\multicolumn{2}{l}{\footnotesize\itshape Tabla \thetable\ (continuación)}\\
\toprule \rowcolor{verdeUTEQ}\color{white}\textbf{Código} & \color{white}\textbf{Enunciado}\\\midrule\endhead
\bottomrule\endfoot
E01 & El paciente podrá reservar una cita eligiendo especialidad, médico, fecha y hora entre los horarios disponibles.\\
E02 & La plataforma enviará al paciente un recordatorio por mensaje de texto 24 horas antes de su cita.\\
E03 & La aplicación debe ser rápida.\\
E04 & El recepcionista podrá registrar la llegada del paciente que se acerca a la ventanilla.\\
E05 & La plataforma se desarrollará en Java 21 con base de datos PostgreSQL, tecnologías que domina el área de TI.\\
E06 & El médico podrá consultar la lista de pacientes del día y marcar cada consulta como atendida.\\
E07 & Los datos de salud de los pacientes deben estar seguros.\\
E08 & La aplicación móvil debe funcionar en teléfonos con Android 10 o superior.\\
E09 & Una cita cancelada con menos de cuatro horas de anticipación se registrará como inasistencia.\\
E10 & La plataforma debe ser fácil de usar para los pacientes adultos mayores.\\
E11 & El proyecto debe entregarse en un plazo máximo de seis meses.\\
E12 & La sesión de un usuario se cerrará automáticamente tras 10 minutos sin actividad.\\
E13 & El quiosco de la sala de espera leerá el número de cédula ingresado en su pantalla táctil e imprimirá un ticket con el número de turno.\\
E14 & La plataforma debe estar siempre disponible.\\
E15 & El director médico podrá generar un reporte mensual de inasistencias por especialidad.\\
E16 & Cada incremento se validará con el director médico al final de cada iteración de dos semanas.\\
E17 & El tratamiento de los datos personales de los pacientes debe cumplir la Ley Orgánica de Protección de Datos Personales del Ecuador.\\
E18 & La plataforma debe integrarse sin problemas con el sistema de facturación actual.\\
E19 & Si el servidor principal falla, la plataforma debe restablecer el servicio en menos de 30 minutos sin perder ninguna cita confirmada.\\
E20 & Todo cambio en el código fuente será revisado por otro integrante del equipo antes de integrarse a la rama principal.\\
\end{longtable}
}

\section{Clasificación en tres dimensiones y justificación}
En la columna ISO se identifica la característica y la subcaracterística solo en los requisitos de calidad; un guion significa que no corresponde asignar esta taxonomía como clasificación principal. En E07 se seleccionan dos subcaracterísticas porque ``seguros'' abarca confidencialidad e integridad. E19 ya contiene una métrica de recuperación, por lo que se clasifica como calidad, aunque no está entre los cinco enunciados vagos que deben reescribirse.

{\fontsize{8.6}{10.3}\selectfont
\setlength\tabcolsep{2.6pt}
\begin{longtable}{@{}P{0.79cm}P{1.95cm}P{1.48cm}P{1.58cm}P{4.05cm}P{6.15cm}@{}}
\caption{Clasificación argumentada de los 20 enunciados}\label{tab:clasificacion}\\
\toprule \rowcolor{verdeUTEQ}
\color{white}\textbf{ID} & \color{white}\textbf{Naturaleza} & \color{white}\textbf{P/P} & \color{white}\textbf{Nivel} & \color{white}\textbf{ISO/IEC 25010:2023} & \color{white}\textbf{Justificación}\\\midrule\endfirsthead
\multicolumn{6}{l}{\footnotesize\itshape Tabla \thetable\ (continuación)}\\
\toprule \rowcolor{verdeUTEQ}
\color{white}\textbf{ID} & \color{white}\textbf{Naturaleza} & \color{white}\textbf{P/P} & \color{white}\textbf{Nivel} & \color{white}\textbf{ISO/IEC 25010:2023} & \color{white}\textbf{Justificación}\\\midrule\endhead
\bottomrule\endfoot
E01 & Funcional & Producto & Software & --- & Permite realizar la reserva y determina los campos de selección que procesa la aplicación.\\
E02 & Funcional & Producto & Sistema & --- & Ordena emitir SMS a través de un canal de comunicación externo, 24 h antes de la cita.\\
\rowcolor{grisTenue} E03 & Calidad & Producto & Software & Eficiencia de desempeño / comportamiento temporal & ``Rápida'' alude al tiempo de respuesta de las operaciones digitales, sin dar un valor objetivo.\\
E04 & Funcional & Producto & Software & --- & Define la función para registrar en el sistema la presencia del paciente.\\
E05 & Restricción & Producto & Software & --- & Fija Java 21 y PostgreSQL, limitando las tecnologías de implementación del producto.\\
E06 & Funcional & Producto & Software & --- & Establece dos acciones del médico: consultar su agenda y registrar una atención.\\
\rowcolor{grisTenue} E07 & Calidad & Producto & Sistema & Seguridad / confidencialidad e integridad & La protección de datos sanitarios alcanza aplicaciones, almacenamiento y accesos del conjunto.\\
E08 & Restricción & Producto & Software & --- & Impone una plataforma de ejecución mínima: Android 10 o superior.\\
E09 & Funcional & Producto & Software & --- & Codifica una regla de negocio que cambia el estado de una cita según el plazo de cancelación.\\
\rowcolor{grisTenue} E10 & Calidad & Producto & Software & Capacidad de interacción / inclusividad y facilidad de aprendizaje & Exige que las personas mayores puedan interactuar y aprender el flujo de reserva.\\
E11 & Restricción & Proceso & No aplica & --- & Fija la duración máxima del proyecto, no una propiedad del sistema terminado.\\
E12 & Funcional & Producto & Software & --- & Especifica una acción automática comprobable: cerrar la sesión a los 10 min de inactividad.\\
E13 & Funcional & Producto & Sistema & --- & Combina software, pantalla táctil e impresora térmica para entregar un turno físico.\\
\rowcolor{grisTenue} E14 & Calidad & Producto & Sistema & Fiabilidad / disponibilidad & ``Siempre disponible'' expresa continuidad del servicio, sin porcentaje ni período definidos.\\
E15 & Funcional & Producto & Software & --- & Define un reporte que puede generar el director, filtrado por mes y especialidad.\\
E16 & Restricción & Proceso & No aplica & --- & Establece validación y periodicidad de iteraciones durante el desarrollo.\\
E17 & Restricción & Producto & Sistema & --- & Impone conformidad legal al tratamiento de información personal en toda la solución.\\
\rowcolor{grisTenue} E18 & Calidad & Producto & Sistema & Compatibilidad / interoperabilidad & Se refiere al intercambio correcto de cobros con otro sistema mediante servicios web.\\
E19 & Calidad & Producto & Sistema & Fiabilidad / recuperabilidad & Define recuperación ante falla de servidor: menos de 30 min y cero citas confirmadas perdidas.\\
E20 & Restricción & Proceso & No aplica & --- & Exige revisión humana de cambios de código antes de integrarlos; regula el proceso de software.\\
\end{longtable}
}
\noindent\textbf{Aclaraciones relevantes.} E05 y E08 son restricciones de tecnología/entorno, aunque inciden en atributos de calidad. E12 es funcional porque prescribe el comportamiento exacto de expiración de sesión; la seguridad es su motivación, no su clasificación principal. E17 es una restricción normativa, no una lista implícita de controles técnicos. Para E11, E16 y E20, el nivel sistema/software se deja como ``no aplica'' al tratarse de procesos del proyecto; en E20 el proceso específico es el desarrollo de software. Estas decisiones deberán contrastarse con el compañero.

\section{Reformulación cuantificada de los cinco enunciados vagos}
Los umbrales de esta sección son propuestas razonadas para el ejercicio, basadas en los 300 usuarios concurrentes, el horario de atención, la proporción de adultos mayores, la sensibilidad de la información y la integración existente. No figuran cuantificados en el caso original ni se ha obtenido aprobación del director médico. Antes de convertirlos en líneas base contractuales deben ser validados con los interesados \cite{guia2026,iso29148}.

\subsection*{RC-03 (origen E03): rendimiento de la reserva}
\textbf{Redacción verificable.} Con 300 sesiones de usuario simultáneas en el ambiente de prueba representativo y consultas válidas al catálogo de médicos y horarios, al menos el 95\,\% de las solicitudes de disponibilidad deberán mostrar una respuesta completa en 3 segundos o menos, medidos desde el envío de la solicitud hasta su presentación en el cliente.
\begin{itemize}
\item \textbf{Característica/subcaracterística:} eficiencia de desempeño / comportamiento temporal.
\item \textbf{Métrica, unidad y umbral:} percentil 95 (P95) del tiempo de respuesta; segundos; \(\mathrm{P95}\leq 3\,\mathrm{s}\).
\item \textbf{Verificación propuesta:} ejecutar al menos 1\,000 solicitudes registradas durante una prueba de carga con 300 sesiones activas; calcular el P95 sobre tiempos válidos. Incluir un dispositivo Android de gama media/baja en la validación de experiencia móvil.
\end{itemize}

\subsection*{RC-07 (origen E07): protección de información de salud}
\textbf{Redacción verificable.} El 100\,\% de las comunicaciones que transfieran datos de salud entre clientes y servicios deberá utilizar TLS 1.2 o superior; el 100\,\% de los registros que contengan datos de salud persistidos deberá encontrarse cifrado en reposo. En un conjunto de 100 pruebas negativas de autorización previamente definidas, el número de lecturas o modificaciones no autorizadas que tengan éxito deberá ser cero.
\begin{itemize}
\item \textbf{Característica/subcaracterística:} seguridad / confidencialidad e integridad.
\item \textbf{Métrica, unidad y umbral:} cobertura de cifrado en tránsito y reposo, en porcentaje (100\,\% cada una); accesos o modificaciones indebidos exitosos, en eventos (\(0/100\)).
\item \textbf{Verificación propuesta:} revisar configuraciones de transporte y almacenamiento, capturas técnicas de conexión sin incluir datos personales reales, y probar la matriz de permisos con identidades sintéticas. Las medidas no sustituyen una evaluación jurídica completa de la LOPDP \cite{lopdp2021}.
\end{itemize}

\subsection*{RC-10 (origen E10): interacción de pacientes mayores}
\textbf{Redacción verificable.} En una prueba de usabilidad con 20 participantes de 60 años o más, al menos 18 de 20 (90\,\%) deberán conseguir reservar una cita disponible, desde la pantalla inicial hasta la confirmación, sin ayuda y en un máximo de 5 minutos por persona.
\begin{itemize}
\item \textbf{Característica/subcaracterística:} capacidad de interacción / inclusividad y facilidad de aprendizaje.
\item \textbf{Métrica, unidad y umbral:} tasa de éxito autónomo, porcentaje (\(\geq90\)\,\%); tiempo por tarea, minutos (\(\leq5\)); muestra objetivo, personas (\(n=20\)).
\item \textbf{Verificación propuesta:} definir una misma tarea y criterio de éxito, cronometrar de forma anónima, registrar errores y calcular éxitos válidos. \emph{Esta prueba todavía no se ha realizado y no se alegan resultados.}
\end{itemize}

\subsection*{RC-14 (origen E14): disponibilidad operativa}
\textbf{Redacción verificable.} La plataforma completa de reservas deberá alcanzar una disponibilidad mensual de al menos 99,9\,\% durante el horario del policlínico (lunes a sábado, 07:00--19:00), medida mediante comprobaciones automáticas cada 60 segundos. Una comprobación será exitosa si el servicio permite consultar disponibilidad en 5 segundos o menos.
\begin{itemize}
\item \textbf{Característica/subcaracterística:} fiabilidad / disponibilidad.
\item \textbf{Métrica, unidad y umbral:} porcentaje de intervalos operativos exitosos por mes (\(\geq99{,}9\)\,\%); intervalo de medición, segundos (60).
\item \textbf{Verificación propuesta:} contar las comprobaciones exitosas y totales del horario definido: \(A=100\times\tfrac{N_{\mathrm{exitosas}}}{N_{\mathrm{totales}}}\). No excluir fallas ni mantenimiento del denominador salvo acuerdo documentado.
\end{itemize}

\subsection*{RC-18 (origen E18): interoperabilidad con facturación}
\textbf{Redacción verificable.} En una prueba de integración con la interfaz web documentada del sistema de facturación existente, la plataforma deberá registrar correctamente al menos 198 de 200 solicitudes válidas de cobro (99\,\%) y no generar ningún cobro duplicado cuando se repita una solicitud con el mismo identificador de operación.
\begin{itemize}
\item \textbf{Característica/subcaracterística:} compatibilidad / interoperabilidad.
\item \textbf{Métrica, unidad y umbral:} transacciones válidas registradas con confirmación, porcentaje (\(\geq99\)\,\%, sobre 200 intentos); duplicados, cantidad (0).
\item \textbf{Verificación propuesta:} usar la interfaz de pruebas pactada con el proveedor, validar el identificador de cobro y contrastar solicitudes con los registros del sistema emisor. Confirmar previamente que el servicio permite deduplicación o acordar cómo la implementará la nueva plataforma.
\end{itemize}

\begin{center}
\begin{tabularx}{\textwidth}{@{}P{1.43cm}P{3.0cm}P{4.93cm}X@{}}
\toprule\rowcolor{verdeUTEQ}
\color{white}\textbf{Origen} & \color{white}\textbf{Característica} & \color{white}\textbf{Métrica y unidad} & \color{white}\textbf{Umbral}\\\midrule
E03 & Desempeño temporal & P95 del tiempo de respuesta (s) & \(\leq 3\,\mathrm{s}\), con 300 sesiones\\
E07 & Seguridad & Cobertura de cifrado (\%); fallas de autorización (eventos) & 100\% en ambos medios; 0 de 100\\
E10 & Capacidad de interacción & Éxitos autónomos (\%); duración (min) & \(\geq90\)\,\% y \(\leq5\) min\\
E14 & Disponibilidad & Comprobaciones exitosas del mes (\%) & \(\geq99{,}9\)\,\%\\
E18 & Interoperabilidad & Cobros confirmados (\%); duplicados (n) & \(\geq99\)\,\%; 0 duplicados\\
\bottomrule
\end{tabularx}
\end{center}

\section{Comparación con un compañero y resolución de discrepancias}


\textbf{Compañero/a:} Jeanpierre Robinson Espinoza \hfill \textbf{Fecha real:} 5 de noviembre de 2026

{\small
\begin{longtable}{@{}P{0.8cm}P{2.9cm}P{2.9cm}P{2.5cm}P{6.0cm}@{}}
\toprule \rowcolor{verdeUTEQ}
\color{white}\textbf{ID} & \color{white}\textbf{Mi criterio} & \color{white}\textbf{Criterio del compañero} & \color{white}\textbf{Discrepancia} & \color{white}\textbf{Acuerdo final y argumento}\\\midrule\endfirsthead
\toprule \rowcolor{verdeUTEQ}\color{white}\textbf{ID} & \color{white}\textbf{Mi criterio} & \color{white}\textbf{Criterio del compañero} & \color{white}\textbf{Discrepancia} & \color{white}\textbf{Acuerdo final y argumento}\\\midrule\endhead
\bottomrule\endfoot
\ E12  & Funcional, producto, software & Calidad, producto, software & Naturaleza & Se acuerda funcional porque establece una acción concreta: cerrar la sesión después de 10 minutos sin actividad.\\[1.0cm]
\ E13  & Funcional, producto, sistema & Funcional, producto, software & Nivel & Se acuerda sistema porque el requisito incluye pantalla táctil, software e impresora térmica.\\[1.0cm]
\ E17  & Restricción, producto, sistema & Calidad, producto, sistema & Naturaleza & Se acuerda restricción porque impone el cumplimiento de una ley, sin especificar directamente un atributo medible.\\[1.0cm]
\end{longtable}
}




\section{Declaración transparente de asistencia con IA}
Se utilizó ChatGPT (OpenAI, GPT-6) como apoyo sustantivo para estructurar el borrador inicial del documento y preparar el código \LaTeX{}. La identidad del caso y sus 20 enunciados proceden de la guía académica. Se contrastaron las denominaciones de calidad con la edición ISO/IEC 25010:2023 y los datos bibliográficos con las páginas oficiales de las fuentes; asimismo, se preparó el documento para compilación con \texttt{pdflatex} y \texttt{bibtex}. 

\bibliographystyle{IEEEtran}
\bibliography{referencias}
\end{document}
